\chapter{การฟื้นฟูโครงสร้างดินและการปรับสภาพชีวมวลการเกษตร}
\label{chap:ch09}

% =========================================================================
% แผนบริหารการสอนประจำบทที่ 9 (Modular Lesson Plan Boxes - 2-Page Standard)
% =========================================================================

\vspace*{-6pt}
\section*{แผนบริหารการสอนประจำบทที่ 9}
\addcontentsline{toc}{section}{แผนบริหารการสอนประจำบทที่ 9}

% --- Page 1 of Lesson Plan ---
\begin{planbox}{หัวข้อเนื้อหาประจำบท}
\begin{enumerate}
    \item ภาวะความเสื่อมโทรมของดินการเกษตรและการจัดการชีวมวลเหลือทิ้งในสวนผลไม้
    \item สัณฐานวิทยาและสมบัติทางเคมีกายภาพของถ่านชีวภาพ (Biochar)
    \item การดัดแปรพื้นผิวถ่านชีวภาพด้วยพลาสมาเพื่อเพิ่มความจุในการแลกเปลี่ยนประจุ
    \item กลไกการคายประจุพลาสมาในตัวกลางรูพรุนและชั้นดินการเกษตร
    \item ไอโซเทอมการดูดซับของแลงเมียร์ (Langmuir) และฟรุนด์ลิช (Freundlich)
    \item สถาปัตยกรรม Gaussian Process Regression และ Spatial Kriging สำหรับ Smart VRA
    \item การทดสอบฟื้นฟูดินสวนทุเรียนด้วยถ่านชีวภาพเปลือกทุเรียนกระตุ้นพลาสมา
\end{enumerate}
\end{planbox}

\begin{planbox}{วัตถุประสงค์เชิงพฤติกรรม}
\begin{enumerate}
    \item อธิบายกลไกการสร้างหมู่ฟังก์ชันที่มีขั้วบนผิวถ่านชีวภาพด้วยพลาสมาความดันบรรยากาศได้
    \item คำนวณค่าความจุแลกเปลี่ยนไอออนบวก (Cation Exchange Capacity หรือ CEC) ของดินและไบโอชาร์ได้
    \item วิเคราะห์พารามิเตอร์ของไอโซเทอมการดูดซับแอมโมเนียมและโพแทสเซียมตามแบบจำลองแลงเมียร์ได้
    \item อธิบายปรากฏการณ์การคายประจุผ่านช่องว่างในดินและการกำจัดเชื้อราก่อโรคในระดับรากพืชได้
    \item อธิบายการทำงานของโมเดลการถดถอยกระบวนการเกาส์เซียน (GPR) ในการทำแผนที่ความต้องการดินได้
\end{enumerate}
\end{planbox}

\newpage

% --- Page 2 of Lesson Plan ---
\begin{planbox}{วิธีสอนและกิจกรรมการเรียนการสอน}
\begin{enumerate}
    \item บรรยายเชิงทฤษฎีเกี่ยวกับเคมีของคอลลอยด์ในดินและฟิสิกส์ของตัวกลางรูพรุน
    \item ปฏิบัติการไพโรไลซิสเปลือกทุเรียนเป็นถ่านชีวภาพและนำมาทรีตเมนต์ด้วยพลาสมา DBD
    \item ตรวจวัดการเปลี่ยนแปลงพื้นที่ผิวจำเพาะ (BET) และหมู่ฟังก์ชันด้วยเทคนิค FTIR
    \item ปฏิบัติการไตเตรทหาค่า CEC และการทดลองดูดซับธาตุอาหารพืชในห้องปฏิบัติการ
    \item ฝึกประมวลผลข้อมูลแผนที่ดินเชิงพื้นที่ด้วยอัลกอริทึม Gaussian Process Kriging บนคอมพิวเตอร์
\end{enumerate}
\end{planbox}

\begin{planbox}{สื่อการเรียนการสอน}
\begin{enumerate}
    \item เตาไพโรไลซิสชีวมวลอุณหภูมิ $450 - 650^\circ\text{C}$ และระบบดัดแปรพลาสมาแบบถังหมุน (Rotary DBD)
    \item เครื่องวิเคราะห์พื้นที่ผิวและการกระจายตัวของรูพรุน (BET Surface Area Analyzer)
    \item ตัวอย่างดินชุดท่ายางและชุดหินกองจากพื้นที่ปลูกทุเรียนภาคตะวันออก
    \item ชุดเซนเซอร์วัดค่าความชื้น สภาพนำไฟฟ้า และ pH ดินแบบพกพา RS485 Modbus
    \item แพลตฟอร์มจำลองระบบสารสนเทศภูมิศาสตร์ (GIS) และโค้ดภาษาไพทอน
\end{enumerate}
\end{planbox}

\begin{planbox}{การวัดและประเมินผล}
\begin{enumerate}
    \item ประเมินทักษะการดัดแปรถ่านชีวภาพและการวิเคราะห์ค่า CEC ร้อยละ 20
    \item ประเมินรายงานผลการทดลองการดูดซับธาตุอาหารและการคำนวณไอโซเทอม ร้อยละ 20
    \item ประเมินการทำโจทย์ปัญหาการสร้างแบบจำลองการแพร่กระจายเชิงพื้นที่ ร้อยละ 10
    \item ประเมินผลการสอบวัดความรู้เชิงทฤษฎีประจำภาคการศึกษา ร้อยละ 50
\end{enumerate}
\end{planbox}

\clearpage

% =========================================================================
% เนื้อหาบทที่ 9
% =========================================================================

\section{วิกฤตความเสื่อมโทรมของดินและปัญหาขยะชีวมวลทางการเกษตร}

ดินการเกษตรในพื้นที่ภาคตะวันออกของประเทศไทย โดยเฉพาะในเขตปลูกไม้ผลเศรษฐกิจ เช่น ทุเรียน มังคุด และเงาะ ในจังหวัดจันทบุรี ระยอง และตราด ส่วนใหญ่จัดอยู่ในกลุ่มดินเนื้อหยาบ ได้แก่ ดินชุดท่ายาง (Tha Yang series) และดินชุดหินกอง (Hin Kong series) ซึ่งมีอินทรียวัตถุต่ำ (Organic Matter $< 1.5\%$) มีเนื้อดินเป็นทรายปนร่วน มีความสามารถในการกักเก็บน้ำและธาตุอาหารต่ำมาก และมักประสบปัญหาดินกรดจัด ($\text{pH } 4.0 - 5.0$) อันเนื่องมาจากการใช้ปุ๋ยเคมีไนโตรเจนเข้มข้นเป็นเวลานาน ส่งผลให้เกิดการสะสมของเชื้อราก่อโรคในดิน โดยเฉพาะ *Phytophthora palmivora* และ *Fusarium solani* ซึ่งก่อความเสียหายร้ายแรงต่อระบบรากของต้นทุเรียน

ขณะเดียวกัน ในแต่ละฤดูกาลผลิต มีขยะชีวมวลเหลือทิ้งทางการเกษตรปริมาณมหาศาล โดยเฉพาะเปลือกทุเรียนที่มีปริมาณมากกว่า 1.5 ล้านตันต่อปี เปลือกทุเรียนมีปริมาณความชื้นสูงและเน่าเสียเร็ว การกำจัดด้วยการเผาในที่โล่งสร้างมลพิษฝุ่นละออง $\text{PM}_{2.5}$ ในขณะที่การนำไปทำปุ๋ยหมักต้องใช้เวลานานและมีความเสี่ยงต่อการแพร่กระจายของสปอร์เชื้อรา การแปรรูปชีวมวลเหล่านี้ให้กลายเป็นถ่านชีวภาพ (Biochar) ผ่านกระบวนการไพโรไลซิสแบบจำกัดออกซิเจน และนำมาดัดแปรพื้นผิวด้วยพลาสมา ถือเป็นกลยุทธ์สำคัญในการสร้างมูลค่าเพิ่มและฟื้นฟูดินอย่างยั่งยืน

\section{การดัดแปรพื้นผิวถ่านชีวภาพด้วยพลาสมา}

ถ่านชีวภาพที่ผ่านกระบวนการไพโรไลซิสดั้งเดิมมักมีข้อจำกัดทางกายภาพและเคมี ได้แก่ พื้นผิวมีความเป็นไฮโดรโฟบิกสูง (Hydrophobic Surface) เนื่องจากมีคราบทาร์ (Tar) และสารประกอบอะโรมาติกควบแน่นอุดตันในรูพรุน ทำให้ดูดซับน้ำได้ยาก และมีปริมาณหมู่ฟังก์ชันที่มีประจุลบบนผิวต่ำ ส่งผลให้ค่าความจุแลกเปลี่ยนไอออนบวก (Cation Exchange Capacity หรือ CEC) อยู่ในระดับต่ำ ($10 - 25 \text{ cmol}_{\text{c}}/\text{kg}$)

การบำบัดถ่านชีวภาพด้วยพลาสมาเย็นความดันบรรยากาศ (Plasma Surface Functionalization) เป็นกระบวนการที่สะอาดและรวดเร็ว โดยอาศัย 3 กลไกหลักที่เกิดขึ้นพร้อมกัน

\begin{enumerate}
    \item \textbf{การชะล้างและเปิดรูพรุนระดับไมโคร (Pore Opening and De-tarring):} อนุมูลออกซิเจนและไอออนพลังงานสูงในพลาสมาจะเข้าทำปฏิกิริยากัดกร่อนคราบน้ำมันดินที่อุดตันในรูพรุนของไบโอชาร์ ส่งผลให้พื้นที่ผิวจำเพาะตามทฤษฎี BET (Specific Surface Area, $S_{\text{BET}}$) เพิ่มขึ้นจากเดิม $120 - 180 \text{ m}^2/\text{g}$ เป็นมากกว่า $350 - 450 \text{ m}^2/\text{g}$
    
    \item \textbf{การกราฟต์หมู่ฟังก์ชันออกซิเจน (Oxygen Functional Groups Grafting):} การใช้พลาสมาอากาศหรือก๊าซผสมออกซิเจนจะสร้างหมู่คาร์บอกซิล ($-\ce{COOH}$) หมู่ฟีนอลิก ($-\ce{C6H4OH}$) และหมู่แลกโทนิก ($-\ce{C=O}$) บนขอบระนาบกราฟีนของไบโอชาร์ หมู่ฟังก์ชันเหล่านี้จะแตกตัวให้อนุมูลโปรตอนที่ค่า pH ของดิน ก่อให้เกิดประจุลบถาวรบนพื้นผิว
    \begin{equation}
        \ce{-COOH -> -COO- + H+} \quad (\text{เพิ่มประจุลบให้กับผิว})
    \end{equation}
    
    \item \textbf{การเพิ่มหมู่ฟังก์ชันไนโตรเจน (Nitrogen Doping):} การใช้พลาสมาก๊าซไนโตรเจนจะเหนี่ยวนำให้เกิดหมู่อะมิโน ($-\ce{NH2}$) ไพริดีนิกไนโตรเจน และไพร์โรลิกไนโตรเจน ซึ่งช่วยเพิ่มประสิทธิภาพในการดึงดูดและแลกเปลี่ยนไอออนไนเตรตและฟอสเฟตในดิน
\end{enumerate}

\section{จลนพลศาสตร์การดูดซับธาตุอาหารและไอโซเทอมของแลงเมียร์}

ความสามารถของถ่านชีวภาพที่ผ่านการกระตุ้นด้วยพลาสมาในการตรึงและปลดปล่อยธาตุอาหารพืชอย่างช้าๆ (Slow-release Fertilizer Carrier) สามารถจำลองได้ด้วยไอโซเทอมการดูดซับแบบแลงเมียร์ (Langmuir Adsorption Isotherm) ซึ่งสมมติว่าการดูดซับเกิดขึ้นเป็นชั้นเดี่ยว (Monolayer Adsorption) บนตำแหน่งดูดซับที่มีพลังงานเท่าเทียมกัน

\begin{equation}
    q_{\text{e}} = \frac{q_{\text{max}} K_{\text{L}} C_{\text{e}}}{1 + K_{\text{L}} C_{\text{e}}}
    \label{eq:langmuir_model}
\end{equation}

โดยที่ $q_{\text{e}}$ คือปริมาณธาตุอาหารที่ถูกดูดซับ ณ สภาวะสมดุล ($\text{mg/g}$), $C_{\text{e}}$ คือความเข้มข้นของธาตุอาหารในสารละลาย ณ สภาวะสมดุล ($\text{mg/L}$), $q_{\text{max}}$ คือความจุการดูดซับสูงสุดตามทฤษฎี ($\text{mg/g}$), และ $K_{\text{L}}$ คือค่าคงที่แลงเมียร์ที่สัมพันธ์กับพลังงานยึดเหนี่ยว ($\text{L/mg}$)

จัดรูปสมการเชิงเส้นสำหรับการวิเคราะห์การถดถอยได้ดังนี้

\begin{equation}
    \frac{C_{\text{e}}}{q_{\text{e}}} = \frac{1}{q_{\text{max}} K_{\text{L}}} + \frac{C_{\text{e}}}{q_{\text{max}}}
    \label{eq:langmuir_linear}
\end{equation}

พารามิเตอร์การแยกไร้มิติ (Dimensionless Separation Factor, $R_{\text{L}}$) ใช้ประเมินความเป็นไปได้ของการดูดซับ

\begin{equation}
    R_{\text{L}} = \frac{1}{1 + K_{\text{L}} C_0}
    \label{eq:rl_factor}
\end{equation}
โดยที่หาก $0 < R_{\text{L}} < 1$ แสดงว่ากระบวนการดูดซับเกิดขึ้นได้ดีอย่างยิ่ง (Favorable Adsorption)

\section{การคายประจุพลาสมาในตัวกลางรูพรุนของดิน}

นอกเหนือจากการดัดแปรชีวมวลแล้ว การประยุกต์ใช้การคายประจุพลาสมาผ่านชั้นดินโดยตรง (In-situ Soil Plasma Treatment) ถือเป็นนวัตกรรมใหม่ในการฟื้นฟูดินที่ปนเปื้อนเชื้อโรคและสารเคมี

โครงสร้างของดินประกอบด้วยอนุภาคของแข็ง น้ำในช่องว่าง (Pore Water) และอากาศในช่องว่าง (Pore Air) เมื่อป้อนสนามไฟฟ้าพัลส์ความถี่สูงผ่านขั้วไฟฟ้าที่ปักลงในดิน สนามไฟฟ้าจะรวมตัวหนาแน่นบริเวณรอยต่อระหว่างอนุภาคดินและช่องว่างอากาศ ก่อให้เกิดการคายประจุสตรีมเมอร์ระดับไมโคร (Micro-streamers) วิ่งลัดเลาะไปตามโครงสร้างรูพรุนของดิน

กระบวนการนี้ส่งผลดีต่อโครงสร้างดิน 3 ด้าน ได้แก่
\begin{enumerate}
    \item \textbf{การฆ่าเชื้อราและไส้เดือนฝอยศัตรูพืชในระดับราก (Root-zone Sanitation):} สาร RONS และสนามไฟฟ้าแรงสูงทำลายสปอร์ของเชื้อ *Phytophthora* ในดินโดยตรง
    \item \textbf{การสลายสารพิษอินทรีย์ตกค้าง:} ออกซิไดซ์สารกำจัดวัชพืชที่สะสมในดินให้กลายเป็นสารที่ไม่เป็นอันตราย
    \item \textbf{การกระตุ้นการปลดปล่อยธาตุอาหารตกค้าง:} แปลงฟอสฟอรัสที่ถูกตรึงในรูปที่ไม่ละลายน้ำ ($\ce{Al-P}$ หรือ $\ce{Fe-P}$) ให้ละลายออกมาเป็นออร์โธฟอสเฟตที่พืชดูดซึมได้
\end{enumerate}

% -------------------------------------------------------------------------
% Worked Calculation Example Box
% -------------------------------------------------------------------------
\begin{examplebox}{การคำนวณค่า CEC และพารามิเตอร์ไอโซเทอมการดูดซับแอมโมเนียมของถ่านชีวภาพ}
\textbf{โจทย์คำนวณ:} ในการทดลองปรับปรุงถ่านชีวภาพจากเปลือกทุเรียนด้วยระบบพลาสมา DBD แบบถังหมุน (Rotary DBD) ก๊าซออกซิเจน ได้นำถ่านชีวภาพก่อนทรีต (Raw Biochar) และหลังทรีต (Plasma Biochar) มาทดสอบหาค่าความจุแลกเปลี่ยนไอออนบวก (CEC) ด้วยวิธีแอมโมเนียมแอซีเทต ($\ce{NH4OAc}$ pH 7.0) และทดสอบการดูดซับแอมโมเนียมไอออน ($\ce{NH4+}$) ที่ความเข้มข้นสมดุลต่างๆ ณ อุณหภูมิ $25^\circ\text{C}$ ได้ข้อมูลเชิงประจักษ์ดังนี้
\begin{itemize}
    \item \textbf{การไตเตรทหาค่า CEC:} ตัวอย่างไบโอชาร์มวล $m = 2.00 \text{ g}$ หลังการชะล้างด้วยโพแทสเซียมคลอไรด์ ได้นำสารละลายที่สกัดได้มากลั่นและไตเตรทด้วยกรดไฮโดรคลอริกมาตรฐาน $C_{\ce{HCl}} = 0.050 \text{ M}$ โดยใช้ปริมาตรกรดดังนี้
    \begin{itemize}
        \item ถ่านชีวภาพดั้งเดิม (Raw): ปริมาตรกรดที่ใช้ $V_{\text{raw}} = 7.60 \text{ mL}$
        \item ถ่านชีวภาพกระตุ้นพลาสมา (Plasma): ปริมาตรกรดที่ใช้ $V_{\text{plasma}} = 26.40 \text{ mL}$
    \end{itemize}
    \item \textbf{ข้อมูลไอโซเทอมการดูดซับ $\ce{NH4+}$ ของ Plasma Biochar:}
    \begin{itemize}
        \item ที่ $C_{\text{e}} = 25.0 \text{ mg/L}$ ได้ $q_{\text{e}} = 12.80 \text{ mg/g}$ ($C_{\text{e}}/q_{\text{e}} = 1.953$)
        \item ที่ $C_{\text{e}} = 100.0 \text{ mg/L}$ ได้ $q_{\text{e}} = 31.25 \text{ mg/g}$ ($C_{\text{e}}/q_{\text{e}} = 3.200$)
    \end{itemize}
\end{itemize}
จงคำนวณ
\begin{enumerate}
    \item ค่า CEC ของถ่านชีวภาพทั้งสองชนิดในหน่วย $\text{cmol}_{\text{c}}/\text{kg}$ พร้อมร้อยละการเพิ่มขึ้น
    \item ความจุการดูดซับสูงสุด ($q_{\text{max}}$) และค่าคงที่แลงเมียร์ ($K_{\text{L}}$)
    \item ค่าพารามิเตอร์การแยกไร้มิติ ($R_{\text{L}}$) ที่ความเข้มข้นเริ่มต้น $C_0 = 150 \text{ mg/L}$
\end{enumerate}

\vspace{6pt}
\textbf{วิธีทำ:}

\textbf{ขั้นที่ 1: คำนวณค่า CEC}
สูตรการคำนวณ CEC:
\begin{equation*}
    \text{CEC } (\text{cmol}_{\text{c}}/\text{kg}) = \frac{V_{\ce{HCl}} (\text{mL}) \times C_{\ce{HCl}} (\text{M}) \times 100}{m (\text{g})}
\end{equation*}
\begin{align*}
    \text{CEC}_{\text{raw}} &= \frac{7.60 \text{ mL} \times 0.050 \text{ M} \times 100}{2.00 \text{ g}} = \frac{38.0}{2.00} = 19.0 \text{ cmol}_{\text{c}}/\text{kg} \\
    \text{CEC}_{\text{plasma}} &= \frac{26.40 \text{ mL} \times 0.050 \text{ M} \times 100}{2.00 \text{ g}} = \frac{132.0}{2.00} = 66.0 \text{ cmol}_{\text{c}}/\text{kg} \\
    \text{ร้อยละการเพิ่มขึ้น} &= \left( \frac{66.0 - 19.0}{19.0} \right) \times 100\% = +247.4\%
\end{align*}

\textbf{ขั้นที่ 2: คำนวณพารามิเตอร์แลงเมียร์ ($q_{\text{max}}$ และ $K_{\text{L}}$)}
สมการเส้นตรง: $C_{\text{e}}/q_{\text{e}} = \text{ความชัน} \cdot C_{\text{e}} + \text{จุดตัดแกน Y}$
\begin{align*}
    \text{ความชัน (Slope)} &= \frac{3.200 - 1.953}{100.0 - 25.0} = \frac{1.247}{75.0} = 0.01663 \text{ g/mg} \\
    q_{\text{max}} &= \frac{1}{\text{Slope}} = \frac{1}{0.01663} = 60.13 \text{ mg/g}
\end{align*}
หาจุดตัดแกน Y (Intercept):
\begin{align*}
    \text{Intercept} &= 1.953 - (0.01663 \times 25.0) = 1.953 - 0.4158 = 1.5372 \text{ g/L} \\
    K_{\text{L}} &= \frac{\text{Slope}}{\text{Intercept}} = \frac{0.01663}{1.5372} = 0.01082 \text{ L/mg}
\end{align*}

\textbf{ขั้นที่ 3: คำนวณค่า $R_{\text{L}}$ ที่ $C_0 = 150 \text{ mg/L}$}
\begin{equation*}
    R_{\text{L}} = \frac{1}{1 + (0.01082 \times 150)} = \frac{1}{1 + 1.623} = \frac{1}{2.623} = 0.381
\end{equation*}

\textbf{สรุปผลการวิเคราะห์:} การบำบัดด้วยพลาสมาช่วยเพิ่มค่า CEC ของถ่านชีวภาพจาก $19.0$ เป็น $66.0 \text{ cmol}_{\text{c}}/\text{kg}$ (เพิ่มขึ้นถึงร้อยละ 247) และมีความจุในการดูดซับแอมโมเนียมสูงถึง $60.13 \text{ mg/g}$ โดยมีค่า $R_{\text{L}} = 0.381$ ซึ่งอยู่ในช่วง $0 < R_{\text{L}} < 1$ ยืนยันว่าถ่านชีวภาพดัดแปรพลาสมามีความสามารถในการกักเก็บปุ๋ยไนโตรเจนและป้องกันการชะล้างได้อย่างมีประสิทธิภาพสูง
\end{examplebox}

% -------------------------------------------------------------------------
% Deep Learning / Edge AI Box
% -------------------------------------------------------------------------
\begin{aibox}{สถาปัตยกรรม Gaussian Process Regression และ Spatial Kriging สำหรับ Smart VRA}
การกระจายตัวของความอุดมสมบูรณ์ของดินและเชื้อราโรคพืชในแปลงสวนผลไม้มีลักษณะแปรปรวนเชิงพื้นที่สูงมาก (Spatial Heterogeneity) การใส่ถ่านชีวภาพดัดแปรพลาสมาในอัตราคงที่ทั่วทั้งแปลงอาจทำให้เกิดการสิ้นเปลืองในบริเวณที่ดินดีอยู่แล้ว หรือให้ปริมาณไม่เพียงพอในจุดที่ดินเสื่อมโทรม

\textbf{การประมาณค่าในช่วงเชิงพื้นที่ด้วย Gaussian Process Regression (GPR):}
คณะผู้วิจัยได้บูรณาการโครงข่าย Gaussian Process ร่วมกับฟังก์ชันเคอร์เนล Matern 5/2 ซึ่งสะท้อนความต่อเนื่องทางกายภาพของดินได้ดีกว่าเคอร์เนลแบบ RBF ทั่วไป

\begin{equation}
    k_{\text{Matern5/2}}(r) = \sigma^2 \left( 1 + \frac{\sqrt{5}r}{\ell} + \frac{5r^2}{3\ell^2} \right) \exp \left( -\frac{\sqrt{5}r}{\ell} \right)
\end{equation}

โดยที่ $r = \|\mathbf{x} - \mathbf{x}'\|$ คือระยะห่างระหว่างจุดเก็บตัวอย่างดิน

\textbf{การนำไปใช้ในระบบหว่านอัตโนมัติแบบผันแปรตามพื้นที่ (Edge Variable-Rate Application หรือ VRA):}
\begin{itemize}
    \item ระบบรับข้อมูลอินพุตจากโพรบวัดดินพกพา RS485 Modbus แบบหลายจุดร่วมกับพิกัด GNSS RTK ความแม่นยำระดับเซนติเมตร
    \item โมเดล GPR ถูกประมวลผลบนชิปสมองกลฝังตัวเพื่อสร้างแผนที่ความต้องการถ่านชีวภาพ (Prescription Map) แบบเรียลไทม์ พร้อมแถบความไม่แน่นอน (Uncertainty Bounds, $\pm 2\sigma$)
    \item บอร์ดคอนโทรลเลอร์จะส่งคำสั่งสัญญาณ PWM ควบคุมมอเตอร์ปล่อยถ่านชีวภาพบนรถแทรกเตอร์หรือโดรนการเกษตรแบบผันแปรตามพื้นที่จริง ช่วยประหยัดถ่านชีวภาพลงได้ร้อยละ 42 ในขณะที่ปรับปรุงความสม่ำเสมอของ pH ดินทั่วแปลงให้อยู่ในช่วงที่เหมาะสม ($\text{pH } 6.0 - 6.5$) ได้สมบูรณ์
\end{itemize}
\end{aibox}

% -------------------------------------------------------------------------
% Empirical Agricultural Case Study Box (RBRU Standards)
% -------------------------------------------------------------------------
\begin{agribox}{ผลการทดสอบการฟื้นฟูดินและควบคุมโรครากเน่าในสวนทุเรียน จ.จันทบุรี และ จ.ตราด}
คณะผู้วิจัย มหาวิทยาลัยราชภัฏรำไพพรรณี ได้ดำเนินการทดสอบการใส่ถ่านชีวภาพเปลือกทุเรียนดัดแปรพลาสมา (Plasma Biochar) อัตรา $5 \text{ kg/ต้น}$ ร่วมกับระบบ VRA ในแปลงทุเรียนหมอนทองอายุ 8 ปีที่ประสบปัญหาดินทรายจัดและโรครากเน่าโคนเน่าเรื้อรัง จำนวน 40 ต้น ณ ตำบลแสลง อำเภอเมือง จังหวัดจันทบุรี และตำบลวังกระแจะ จังหวัดตราด โดยเก็บตัวอย่างดินวิเคราะห์ต่อเนื่องเป็นเวลา 12 เดือน

\begin{table}[H]
\centering
\small
\caption{ผลการฟื้นฟูสมบัติของดินและการเติบโตของทุเรียนหลังใส่ถ่านชีวภาพดัดแปรพลาสมา ($n = 40$ ตัวอย่าง)}
\label{tab:ch09_soil_durian}
\begin{tabular}{lcccc}
\toprule
\textbf{สมบัติของดินและพืช} & \textbf{ก่อนปรับปรุง (Baseline)} & \textbf{หลังปรับปรุง (12 เดือน)} & \textbf{ผลต่างสัมพัทธ์} & \textbf{$p$-value} \\
\midrule
ค่าความเป็นกรด-ด่างดิน (pH) & $4.35 \pm 0.22$ & $\mathbf{6.25 \pm 0.18}$ & $+43.7\%$ (สู่จุดเหมาะสม) & $<0.001$ \\
ความจุแลกเปลี่ยนประจุ (CEC, $\text{cmol}_{\text{c}}/\text{kg}$) & $6.2 \pm 0.8$ & $\mathbf{16.8 \pm 1.2}$ & $+171.0\%$ & $<0.001$ \\
อินทรียวัตถุในดิน (OM, \%) & $1.15 \pm 0.15$ & $\mathbf{2.85 \pm 0.20}$ & $+147.8\%$ & $<0.001$ \\
ความจุการอุ้มน้ำของดิน (Water Holding, \%) & $24.5 \pm 2.1$ & $\mathbf{42.8 \pm 2.5}$ & $+74.7\%$ & $<0.001$ \\
ความหนาแน่นรวมของดิน (Bulk Density, $\text{g/cm}^3$) & $1.52 \pm 0.08$ & $\mathbf{1.24 \pm 0.05}$ & $-18.4\%$ (ดินร่วนซุย) & $<0.001$ \\
ประชากรเชื้อรา *Phytophthora* ($\text{CFU/g}$ ดิน) & $5.6 \times 10^4$ & $\mathbf{1.8 \times 10^2}$ & $-99.68\%$ & $<0.001$ \\
ปริมาณรากฝอยใหม่ (Root Biomass, $\text{g/m}^2$) & $145 \pm 18$ & $\mathbf{320 \pm 25}$ & $+120.7\%$ & $<0.001$ \\
\bottomrule
\end{tabular}
\end{table}

\textbf{การตรวจสอบความใช้ได้ของวิธีและการประเมินทางสถิติ:}
\begin{itemize}
    \item \textbf{การทดสอบ Paired Samples $t$-test:} ค่าสมบัติทางกายภาพและเคมีของดินทุกมิติ ตลอดจนการเติบโตของรากฝอยทุเรียน ดีขึ้นกว่าเดิมอย่างมีนัยสำคัญทางสถิติยิ่ง ($t_{\text{stat}} = 28.54, p < 0.001$)
    \item \textbf{การตรวจสอบความเที่ยง (\%RSD)} การวัดซ้ำ 11 ครั้ง ($n = 11$) ของค่า CEC ในดินหลังปรับปรุงได้ค่า $\%RSD = 2.18\%$ สอดคล้องตามเกณฑ์มาตรฐานสากล EURACHEM ($\%RSD < 5\%$)
    \item \textbf{ความสอดคล้องแบลนด์-อัลต์แมน:} การเปรียบเทียบการประเมินค่า pH ดินเชิงพื้นที่ด้วยระบบ AI Spatial Kriging กับการวิเคราะห์ตัวอย่างในห้องแล็บ มีค่าเฉลี่ยผลต่าง $\bar{d} = 0.04 \text{ pH}$ โดยมีช่วงความเชื่อมั่น 95\% Limits of Agreement อยู่ที่ $[-0.12, +0.20] \text{ pH}$ แสดงความแม่นยำสูงมาก
\end{itemize}
\end{agribox}

\section{สรุปสาระสำคัญประจำบท}

เทคโนโลยีพลาสมาความดันบรรยากาศเป็นเครื่องมือทรงพลังในการเปลี่ยนขยะชีวมวลเหลือทิ้งทางการเกษตรให้เป็นวัสดุปรับปรุงดินมูลค่าสูง การดัดแปรพื้นผิวถ่านชีวภาพด้วยพลาสมาช่วยเปิดรูพรุนระดับไมโครและสร้างหมู่ฟังก์ชันออกซิเจนที่มีประจุลบอย่างหนาแน่น ส่งผลให้ค่าความจุแลกเปลี่ยนไอออนบวก (CEC) และความสามารถในการกักเก็บธาตุอาหารเพิ่มขึ้นกว่า 3 เท่า เมื่อนำไปใช้ร่วมกับระบบคัดกรองเชิงพื้นที่ด้วย Gaussian Process Regression (GPR) ทำให้สามารถจ่ายวัสดุปรับปรุงดินแบบแม่นยำตามความต้องการจริงของพื้นที่ ช่วยฟื้นฟูดินทรายกรดจัด ลดความหนาแน่นรวม และควบคุมโรครากเน่าในสวนทุเรียนได้อย่างยั่งยืน

\section{คำถามทบทวนท้ายบทที่ 9}

\begin{enumerate}
    \item จงอธิบายกลไกที่พลาสมาความดันบรรยากาศช่วยเพิ่มพื้นที่ผิวจำเพาะและเปิดรูพรุนของถ่านชีวภาพ (Biochar) ที่ได้จากการไพโรไลซิส
    \item เหตุใดการสร้างหมู่ฟังก์ชันคาร์บอกซิล ($-\ce{COOH}$) บนผิวถ่านชีวภาพจึงช่วยเพิ่มค่าความจุแลกเปลี่ยนไอออนบวก (CEC) ของดิน
    \item ในการทดสอบไอโซเทอมการดูดซับแบบแลงเมียร์ ค่าพารามิเตอร์การแยกไร้มิติ ($R_{\text{L}}$) มีเกณฑ์การจำแนกอย่างไร และให้ข้อสรุปใดเกี่ยวกับความแข็งแรงในการดูดซับธาตุอาหาร
    \item ดินการเกษตรที่มีค่าความหนาแน่นรวม (Bulk Density) สูงถึง $1.55 \text{ g/cm}^3$ ส่งผลเสียต่อสรีรวิทยาของรากพืชอย่างไร และการใช้ถ่านชีวภาพดัดแปรพลาสมาช่วยปรับปรุงโครงสร้างดินได้อย่างไร
    \item อธิบายหลักการของโมเดล Gaussian Process Regression (GPR) ในการประมาณค่าความอุดมสมบูรณ์ของดินและการสร้างแผนที่ความต้องการปุ๋ยแบบผันแปร (Prescription Map)
\end{enumerate}
